\title{Group Sequence Policy Optimization}

\begin{abstract}
We present Group Sequence Policy Optimization (GSPO), a reinforcement learning algorithm crafted to train large language models with enhanced stability, efficiency, and effectiveness. Differing from earlier approaches that use importance ratios at the token level, GSPO operates by establishing the importance ratio on the entire sequence's probability, applying clipping, reward allocation, and optimization at the sequence scale. Our findings reveal that GSPO outperforms the GRPO algorithm both in terms of sample efficiency and overall performance, consistently stabilizes RL training for Mixture-of-Experts (MoE) architectures, and offers opportunities to streamline RL system design. These advances have played a crucial role in the outstanding capabilities observed in recent Qwen3 models.
\end{abstract}

\section{Introduction}
\label{sec:intro}

Reinforcement learning (RL) represents a fundamental approach to advancing language models, as evidenced by previous studies \citep{o1,dpsk-r1,qwq32b,qwen3}. Employing large-scale RL enables these models to address complex tasks—including high-level mathematics and programming—by fostering more elaborate and prolonged reasoning.

For RL to be scaled effectively with increasing computational resources, it is critical that training remains stable and resilient throughout the process. Despite this requirement, leading RL methods such as GRPO \citep{grpo} are prone to pronounced instability when applied to very large language models, frequently causing catastrophic, unrecoverable collapses \citep{qwen3, minimax-m1}. Such instability acts as a significant barrier to further enhancing the capabilities of language models via ongoing RL-based training.

This work reveals that GRPO's instability originates from the improper implementation and subsequent invalidation of importance sampling weights within its algorithm structure. This results in elevated training variance, which accumulates as responses grow longer and becomes even more severe due to the clipping process, eventually causing the model to collapse.

To overcome these fundamental issues, we introduce \textbf{Group Sequence Policy Optimization (GSPO)}, a novel reinforcement learning method for training large language models. The principal advancement offered by GSPO is its rigorous formulation of the importance ratio rooted in sequence likelihood \citep{click}, adhering closely to the foundational concept of importance sampling. Furthermore, GSPO derives normalized rewards by calculating the advantages across several generated responses for a given query, thereby securing the coherence between reward assignment at the sequence level and the underlying optimization objective.

The results from our empirical assessment highlight GSPO's pronounced advantages over GRPO regarding training stability, efficiency, and overall performance. Notably, GSPO intrinsically addresses the stability issues associated with reinforcement learning (RL) in extensive Mixture-of-Experts (MoE) architectures, obviating the necessity for intricate stabilization techniques and offering a pathway to streamlined RL infrastructure. These strengths of GSPO have played a key role in the marked enhancements seen in the most recent Qwen3 models. We anticipate that GSPO will serve as a dependable and scalable algorithmic base, fostering the ongoing progress of large-scale RL training with language models.

\section{Preliminaries}

\paragraph{Notation}
Throughout this work, we represent an autoregressive language model with parameters $\theta$ as a policy, $\pi_\theta$. Here, $x$ corresponds to a query, while $\mathcal{D}$ denotes the set containing all queries. The probability assigned to a response $y$ given query $x$ by the policy $\pi_\theta$ is written as $\pi_\theta (y | x)=\prod_{t=1}^{|y|} \pi_\theta (y_t | x, y_{<t} )$, where $|y|$ specifies the token count in response $y$. Each query-response pair $(x, y)$ may be evaluated by a verifier $r$, which assigns a reward $r(x, y)$ within the interval $[0, 1]$.

\paragraph{Proximal Policy Optimization (PPO)} Using samples generated from the old policy $\pi_{\theta_\text{old}}$, PPO \citep{ppo} constrains the policy update within a proximal region of the old policy through the clipping mechanism. Specifically, PPO employs the following objective for policy optimization (we omit the KL regularization term hereinafter for brevity, as it is not the focus of this paper):

\begin{align}
\mathcal{J}_\text{PPO}(\theta) = \mathbb{E}_{ x \sim \mathcal{D},\, y \sim \pi_{\theta_\text{old}}( \cdot | x) }
\left[ \frac{1}{|y|} \sum_{t=1}^{|y|} 
\min \left( w_{t}(\theta) \widehat{A}_{t},  \, \mathrm{clip} \left( w_{t}(\theta), 1 - {\varepsilon}, 1 + {\varepsilon}\right) \widehat{A}_{t} \right)
\right],
\end{align}

In this context, the token importance ratio $w_{t}(\theta)$ is given by $
w_{t}(\theta) = \frac{ \pi_{\theta} (y_{t} | x, y_{<t}) }{ \pi_{\theta_\text{old}} (y_{t} | x,y_{<t})} $, where the advantage $\widehat{A}_{t}$ for token $y_t$ is computed through a separate value model, and $\varepsilon$ denotes the clipping threshold applied to these importance ratios.

A fundamental obstacle encountered with PPO during practical implementation is its strong dependence on the value model. Typically, the value model is architecturally comparable to the policy model, which results in significant demands on both memory and computational resources. Moreover, the algorithm’s performance is contingent upon the accuracy of the estimated value function. Developing a dependable value model is difficult, and the problem is compounded when attempting to scale the model to accommodate longer outputs or increasingly intricate tasks.

\paragraph{Group Relative Policy Optimization (GRPO)} GRPO \citep{grpo} bypasses the need for the value model by computing the relative advantage of each response within a group of responses to the same query. Specifically, GRPO optimizes the following objective:

\begin{align}
\label{equ:grpo}
\mathcal{J}_\text{GRPO}(\theta) = \mathbb{E}_{ x \sim \mathcal{D},\, \{y_i\}_{i=1}^G \sim \pi_{\theta_\text{old}}( \cdot | x) }
\left[ \frac{1}{G} \sum_{i=1}^{G} \frac{1}{|y_i|} \sum_{t=1}^{|y_i|} 
\min \left( w_{i,t}(\theta) \widehat{A}_{i,t},  \, \mathrm{clip} \left( w_{i,t}(\theta), 1 - {\varepsilon}, 1 + {\varepsilon}\right) \widehat{A}_{i,t} \right)
\right],
\end{align}

where $G$ is the number of generated responses to each query $x$ (i.e., the group size), and the importance ratio $w_{i,t}(\theta)$ and advantage $\widehat{A}_{i,t}$ of token $y_{i,t}$ are:

\begin{align}
    w_{i,t}(\theta)=\frac{ \pi_{\theta} (y_{i,t} | x, y_{i,<t}) }{ \pi_{\theta_\text{old}} (y_{i,t} | x,y_{i,<t})},\quad\ 
    \widehat{A}_{i,t} = \widehat{A}_{i} = \frac{r(x, y_i) - \mathrm{mean} \left( \{ r(x, y_i) \}_{i=1}^G \right) }{ \mathrm{std} \left( \{ r(x, y_i) \}_{i=1}^G \right) },
\end{align}

respectively, where all the tokens in $y_i$ share the same advantage as $\widehat{A}_{i}$.

\section{Motivation}
\label{sec:motivation}

As models increase in size and sparsity—such as with Mixture-of-Experts architectures—and as the length of generated responses grows, large rollout batch sizes are required to fully leverage hardware during RL training. To enhance sample efficiency, it is common to subdivide these extensive batches of rollout data into several smaller mini-batches for the purpose of gradient computation. However, this approach gives rise to an off-policy scenario, since the sampled responses $y$ originate from a previous policy $\pi_{\theta_\text{old}}$ rather than the current target policy $\pi_\theta$. This situation underlies the rationale for employing clipping techniques in PPO and GRPO, which serve to restrict the influence of samples that deviate too far from the present policy on the gradient estimation process.

While mechanisms like clipping aim to manage this off-policy discrepancy, we identify a more fundamental issue in GRPO: \textit{its objective is ill-posed}. This problem becomes particularly acute when training large models on long-response tasks, leading to catastrophic model collapse. The ill-posed nature of the GRPO objective stems from a misapplication of importance sampling weights. The principle of importance sampling is to estimate the expectation of a function $f$ under a target distribution $\pi_\text{tar}$ by re-weighting samples drawn from a behavior distribution $\pi_\text{beh}$:

\begin{align}
\mathbb{E}_{z \sim \pi_\text{tar}} \left[ f(z) \right] = \mathbb{E}_{z \sim \pi_\text{beh}} \left[ \frac{ \pi_\text{tar} (z) }{ \pi_\text{beh} (z)} \, f(z) \right].
\end{align}

A key aspect of this approach is the need to average across numerous samples ($N \gg 1$) drawn from the behavior policy $\pi_\text{beh}$ for the importance weight $\frac{ \pi_\text{tar} (z) }{ \pi_\text{beh} (z)}$ to accurately address discrepancies between the behavior and target distributions.

On the other hand, GRPO employs the importance weighting $\frac{ \pi_{\theta} (y_{i,t} | x, y_{i,<t}) }{ \pi_{\theta_\text{old}} (y_{i,t} | x,y_{i,<t})}$ at every individual token step $t$. Because this weight is computed from just one token $y_{i,t}$—a sample of the next-token distribution $\pi_{\theta_\text{old}}( \cdot | x, y_{i, <t})$—it fails to achieve the intended adjustment for off-policy data. Instead, it generates considerable variance in the stochastic gradient estimates, especially for longer sequences, with this issue exacerbated by the presence of a clipping mechanism. In our experiments, such training instability has resulted in catastrophic model collapse that cannot be remedied. Once this collapse occurs, no combination of interventions—such as rolling back to an earlier checkpoint, extensive hyperparameter searches (for example, modifying clipping bounds), increasing sequence length, or adjusting reinforcement learning prompts—has managed to restore meaningful training progress.

These findings point to a foundational flaw in how GRPO formulates its objective. The inefficacy of token-level importance weights underscores a crucial lesson: \textbf{the granularity of the optimization objective must align with the granularity of the reward signal}. Because rewards are issued for full sequences, any attempt to apply off-policy corrections at the level of individual tokens is fundamentally misaligned. Consequently, this insight encourages us to abandon token-based objectives in favor of an approach that applies importance weighting and carries out optimization directly at the \textit{sequence level}.

\section{Algorithm}

\subsection{GSPO: Group Sequence Policy Optimization}

Although the token-level importance weight $\frac{ \pi_{\theta} (y_{i,t} | x, y_{i,<t}) }{ \pi_{\theta_\text{old}} (y_{i,t} | x,y_{i,<t})}$ presents issues within GRPO, our analysis indicates that the \textit{sequence-level} importance weight $\frac{ \pi_{\theta} (y | x) }{ \pi_{\theta_\text{old}} (y | x)}$ carries an unambiguous theoretical interpretation in language generation. Specifically, it quantifies the extent to which a response $y$—sampled from $\pi_{\theta_\text{old}} (\cdot | x)$—differs from those likely under $\pi_{\theta} (\cdot | x)$, thus directly corresponding to the sequence-level reward and effectively representing the basis for the clipping mechanism.

Based on this straightforward observation, we propose the \textbf{Group Sequence Policy Optimization (GSPO)} algorithm. GSPO employs the following sequence-level optimization objective:

\begin{align}
\mathcal{J}_\text{GSPO} (\theta) =
\mathbb{E}_{ x \sim \mathcal{D},\, \{y_i\}_{i=1}^G \sim \pi_{\theta_\text{old}}( \cdot | x) }
\left[ 
\frac{1}{G} \sum_{i=1}^{G}
\min \left( s_{i}(\theta)  \widehat{A}_{i},  \, \mathrm{clip} \left( s_{i}(\theta), 1 - {\varepsilon}, 1 + {\varepsilon} \right) \widehat{A}_{i} \right) 
\right],
\label{equ:gspo}
\end{align}

where we adopt the group-based advantage estimation:

\begin{align}
\widehat{A}_{i} = \frac{r(x, y_i) - \mathrm{mean} \left( \{ r(x, y_i) \}_{i=1}^G \right) }{ \mathrm{std} \left( \{ r(x, y_i) \}_{i=1}^G \right) },
\end{align}

and define the importance ratio $s_{i}(\theta)$ based on sequence likelihood \citep{click}:

\begin{align}
s_{i}(\theta) = \left( \frac{ \pi_{\theta} (y_i | x) }{ \pi_{\theta_\text{old}} (y_i | x)} \right)^{\frac{1}{|y_i|}}
=
\exp \left( \frac{1}{|y_i|} \sum_{t=1}^{|y_i|} \log \frac{ \pi_{\theta} (y_{i,t} | x, y_{i,<t}) }{ \pi_{\theta_\text{old}} (y_{i,t} | x,y_{i,<t})} \right).
\end{align}

Consequently, GSPO implements clipping at the level of entire responses, as opposed to individual tokens, to eliminate highly ``off-policy'' samples during gradient computation. This approach aligns with the reward structure and the optimization process, both of which are applied at the sequence level. It is important to highlight that length normalization is employed in $s_{i}(\theta)$, which serves to mitigate variance and ensures that $s_{i}(\theta)$ stays within a consistent numerical interval. Without such normalization, minor changes in the likelihood of certain tokens could lead to substantial volatility in the overall sequence-level importance ratio, necessitating distinct clipping bounds for responses of varying lengths. Furthermore, it is worth mentioning that the typical clipping intervals in GSPO differ considerably from those used in prior methods such as GRPO, owing to the differing ways importance ratios are defined.

\subsection{Gradient Analysis}
\label{subsec:gradient}

We can derive the gradient of the GSPO objective as follows (clipping is omitted for brevity):

\begin{align}
\nabla_{\theta} \mathcal{J}_\text{GSPO} (\theta)
=&\ 
\nabla_{\theta} \mathbb{E}_{ x \sim \mathcal{D},\, \{y_i\}_{i=1}^G \sim \pi_{\theta_\text{old}}( \cdot | x) }
\left[ 
\frac{1}{G} \sum_{i=1}^{G} s_{i}(\theta) \widehat{A}_{i}
\right] \\
= &\ 
\mathbb{E}_{ x \sim \mathcal{D},\, \{y_i\}_{i=1}^G \sim \pi_{\theta_\text{old}}( \cdot | x) }
\left[ 
\frac{1}{G} \sum_{i=1}^{G}
s_{i}(\theta) \widehat{A}_{i}
\cdot \nabla_{\theta} \log s_{i}(\theta)
\right] \\
=&\ 
\mathbb{E}_{ x \sim \mathcal{D},\, \{y_i\}_{i=1}^G \sim \pi_{\theta_\text{old}}( \cdot | x) }
\left[ 
\frac{1}{G} \sum_{i=1}^{G} 
\left( \frac{ \pi_{\theta} (y_i | x) }{ \pi_{\theta_\text{old}} (y_i | x)} \right)^{\frac{1}{|y_i|}} \widehat{A}_{i} 
\cdot \frac{1}{|y_i|} \sum_{t=1}^{|y_i|} \nabla_{\theta} \log \pi_{\theta} (y_{i,t} | x, y_{i,<t}) 
\right].
\label{equ:gspo_grad}
\end{align}

For comparison, the gradient of the GRPO objective is as follows (note that $\widehat{A}_{i,t} = \widehat{A}_{i}$):

\begin{align}
\nabla_{\theta} \mathcal{J}_\text{GRPO} (\theta) 
=&\ 
\nabla_{\theta} \mathbb{E}_{ x \sim \mathcal{D},\, \{y_i\}_{i=1}^G \sim \pi_{\theta_\text{old}}( \cdot | x) }
\left[ \frac{1}{G} \sum_{i=1}^{G} \frac{1}{|y_i|} \sum_{t=1}^{|y_i|} 
w_{i,t}(\theta) \widehat{A}_{i,t}
\right] \\
=&\
\mathbb{E}_{ x \sim \mathcal{D},\, \{y_i\}_{i=1}^G \sim \pi_{\theta_\text{old}}( \cdot | x) }
\left[ \frac{1}{G} \sum_{i=1}^{G} \widehat{A}_{i} 
\cdot \frac{1}{|y_i|} \sum_{t=1}^{|y_i|} 
\frac{ \pi_{\theta} (y_{i,t} | x, y_{i,<t}) }{ \pi_{\theta_\text{old}} (y_{i,t} | x,y_{i,<t})} 
\nabla_{\theta} \log \pi_{\theta} (y_{i,t} | x, y_{i,<t})  
\right].
\end{align}

Thus, the primary difference between GSPO and GRPO concerns \textit{the manner in which they assign weights to the gradients of token log likelihoods}. In GRPO, each token is assigned an “importance weight” given by $\frac{ \pi_{\theta} (y_{i,t} | x, y_{i,<t}) }{ \pi_{\theta_\text{old}} (y_{i,t} | x,y_{i,<t})}$, which introduces non-uniform weighting. These weights can fall within $(0, 1+\varepsilon]$ when $\widehat{A}_{i} >0$, or within $[1-\varepsilon, +\infty)$ when $\widehat{A}_{i} < 0$, and because the disparities among these weights are not minor, their effects can aggregate during optimization and potentially produce unstable outcomes. Conversely, GSPO applies uniform weighting across all tokens in a given response, thereby removing this inherent source of instability found in GRPO.

\subsection{GSPO-token: A Token-level Objective Variant}

In scenarios like multi-turn RL, we may desire a finer-grained advantage adjustment than the sequence level. To this end, we introduce a token-level objective variant of GSPO, namely \textbf{GSPO-token}, to allow token-wise advantage customization:

\begin{align}
\mathcal{J}_\text{GSPO-token}(\theta)
=
\mathbb{E}_{ x \sim \mathcal{D},\, \{y_i\}_{i=1}^G \sim \pi_{\theta_\text{old}}( \cdot | x) }
\left[ 
\frac{1}{G} \sum_{i=1}^{G} \frac{1}{|y_i|} \sum_{t=1}^{|y_i|} 
\min \left( s_{i,t}(\theta) \widehat{A}_{i,t},  \, \mathrm{clip} \left( s_{i,t}(\theta), 1 - {\varepsilon}, 1 + {\varepsilon} \right) \widehat{A}_{i,t} \right) 
\right],
\label{equ:gspo-token}
\end{align}

where

\begin{align}
s_{i,t}(\theta) 
= \mathrm{sg} \left[ s_{i}(\theta) \right]  \cdot \frac{ \pi_{\theta} (y_{i,t} | x, y_{i,<t}) }{ \mathrm{sg} \left[ \pi_{\theta} (y_{i,t} | x, y_{i,<t}) \right] },
\end{align}

and $\mathrm{sg}[\cdot]$ denotes only taking the numerical value but stopping the gradient, corresponding to the \texttt{detach} operation in PyTorch. The gradient of GSPO-token can be derived as:

\begin{align}
\nabla_{\theta} \mathcal{J}_\text{GSPO-token}(\theta)
=&\ 
\nabla_{\theta} \mathbb{E}_{ x \sim \mathcal{D},\, \{y_i\}_{i=1}^G \sim \pi_{\theta_\text{old}}( \cdot | x) }
\left[ 
\frac{1}{G} \sum_{i=1}^{G}
\frac{1}{|y_i|} \sum_{t=1}^{|y_i|} 
s_{i,t}(\theta) \widehat{A}_{i,t}
\right] \\
=&\ 
\mathbb{E}_{ x \sim \mathcal{D},\, \{y_i\}_{i=1}^G \sim \pi_{\theta_\text{old}}( \cdot | x) }
\left[ 
\frac{1}{G} \sum_{i=1}^{G} s_i (\theta) \cdot \frac{1}{|y_i|} \sum_{t=1}^{|y_i|} 
\widehat{A}_{i,t} \frac{ \nabla_{\theta} \pi_{\theta} (y_{i,t} | x, y_{i,<t}) }{ \pi_{\theta} (y_{i,t} | x, y_{i,<t}) }
\right] \\
=&\ 
\mathbb{E}_{ x \sim \mathcal{D},\, \{y_i\}_{i=1}^G \sim \pi_{\theta_\text{old}}( \cdot | x) }
\left[ 
\frac{1}{G} \sum_{i=1}^{G}  \left( \frac{ \pi_{\theta} (y_i | x) }{ \pi_{\theta_\text{old}} (y_i | x)} \right)^{\frac{1}{|y_i|}}
\cdot \frac{1}{|y_i|} \sum_{t=1}^{|y_i|}
\widehat{A}_{i,t} \nabla_{\theta} \log \pi_{\theta} (y_{i,t} | x, y_{i,<t})  
\right].
\label{equ:gspo-token_grad}
\end{align}

Observe that the expression $\frac{ \pi_{\theta} (y_{i,t} | x, y_{i,<t}) }{ \mathrm{sg} \left[ \pi_{\theta} (y_{i,t} | x, y_{i,<t}) \right] }$ evaluates to 1, thereby making $s_{i,t}(\theta)$ numerically identical to $s_{i}(\theta)$. If we compare Equation~\eqref{equ:gspo} and Equation~\eqref{equ:gspo-token}, as well as Equation~\eqref{equ:gspo_grad} with Equation~\eqref{equ:gspo-token_grad}, it is evident that GSPO-token and GSPO yield equivalent numerical results in terms of the objective function, clipping condition, and theoretical gradient, provided that every token's advantage in the response $y_i$ is set to the same value, namely $\widehat{A}_{i,t} = \widehat{A}_{i}$. Nevertheless, GSPO-token provides the added benefit of allowing for token-level advantage customization, enhancing its flexibility.

\section{Experiments and Discussion}

\subsection{Empirical Results}

Our experiments utilize a cold-start model that has been fine-tuned from Qwen3-30B-A3B-Base. For evaluation, we provide both the progression of training rewards and the corresponding model performance on several benchmarks: AIME'24 (average Pass@1 across 32 sampled attempts), LiveCodeBench (from October 2024 to February 2025, average Pass@1 across 8 samples), and CodeForces (using Elo Rating). Throughout the reinforcement learning phase, we subdivide each rollout batch into four separate mini-batches for the purpose of gradient updates. In our GSPO implementation, the clipping intervals in Equation~\eqref{equ:gspo} are specified as 3e-4 for the lower bound and 4e-4 for the upper bound. GRPO serves as our baseline comparison, with its clipping parameters in Equation~\eqref{equ:grpo} established at 0.2 (left) and 0.27 (right); these values have been extensively fine-tuned to facilitate an equitable evaluation. Importantly, GRPO is dependent on the Routing Replay training technique to achieve proper convergence within MoE RL, a requirement which will be explored further in \S~\ref{subsec:routing_replay}. In contrast, \textit{GSPO no longer requires this additional training strategy}.

As illustrated in Figure~\ref{fig:results}, GSPO maintains a stable progression during training. Notably, \textbf{GSPO enables sustained improvements in performance by augmenting training compute, periodically refreshing the query set, and allowing for increased generation lengths}. Additionally, GSPO exhibits higher training efficiency compared to GRPO, reaching superior training accuracy and benchmark outcomes with equivalent amounts of compute and queries consumed. The application of GSPO to RL training for the latest Qwen3 models further confirms GSPO’s effectiveness in harnessing RL scaling capabilities within large language models.

\subsection{Curious Observation on Clipping Fractions}

One of the primary differences between GSPO and GRPO lies in the scope of their clipping procedures: GSPO operates by clipping entire sequences, whereas GRPO applies clipping at the individual token level. As illustrated in Figure~\ref{fig:clipping}, this methodological divergence yields a dramatic contrast—by about two orders of magnitude—in the proportions of clipped tokens between the two approaches, a gap that persists regardless of how the clipping intervals are set. Remarkably, although GSPO discards a substantially larger number of tokens and therefore utilizes fewer for both training and gradient estimation, it nonetheless achieves greater training efficiency than GRPO. This somewhat surprising outcome—where the exclusion of a higher fraction of tokens corresponds to better training performance—strongly suggests that the gradient estimates derived at the token level in GRPO are characterized by considerable noise and thus provide a suboptimal basis for learning. Conversely, GSPO’s reliance on sequence-level information furnishes more consistent and effective training signals.

\subsection{Benefit of GSPO for MoE Training}
\label{subsec:routing_replay}

\paragraph{Background}
Unlike the dense models typically used in RL training, MoE models, due to their sparsely activated architectures, present distinct challenges related to training stability. Our investigations revealed that, with the GRPO algorithm, MoE models can suffer from pronounced \textit{expert-activation volatility}, which disrupts the convergence of RL optimization. Specifically, after performing one or several steps of gradient update, the set of experts selected for producing the identical output may shift considerably. For instance, in the case of the 48-layer Qwen3-30B-A3B-Base model, each RL gradient update leads to approximately 10\% of the experts activated for the same rollout differing under the updated policy $\pi_{\theta}$ compared to the previous policy $\pi_{\theta_\text{old}}$.

This behavior is increasingly evident as the depth of the MoE architecture grows. Consequently, it causes significant oscillations in the token-level importance weights $w_{i,t}(\theta) = \frac{ \pi_{\theta} (y_{i,t} | x, y_{i,<t}) }{ \pi_{\theta_\text{old}} (y_{i,t} | x,y_{i,<t})}$, rendering them unreliable, as elaborated in \S~\ref{sec:motivation} and~\ref{subsec:gradient}. These dramatic fluctuations subsequently interfere with effective RL training convergence.

\paragraph{Our Previous Approach} In our earlier work, we addressed this issue using the \textbf{Routing Replay} training scheme. This method involves storing the expert selection—specifically, the activated experts—determined by $\pi_{\theta_\text{old}}$, and ensuring that the identical routing configurations are employed in $\pi_{\theta}$ when calculating the importance weights $w_{i,t}(\theta) = \frac{ \pi_{\theta} (y_{i,t} | x, y_{i,<t}) }{ \pi_{\theta_\text{old}} (y_{i,t} | x,y_{i,<t})}$. As a result, both $\pi_{\theta} (y_{i,t} | x, y_{i,<t})$ and $\pi_{\theta_\text{old}} (y_{i,t} | x,y_{i,<t})$ for each token $y_{i,t}$ rely on an identical set of activated network components, which stabilizes the token-wise importance weights and promotes consistent optimization over the same expert pathways throughout gradient updates. As depicted in Figure~\ref{fig:routing_replay}, Routing Replay is a pivotal component for achieving stable and effective convergence during GRPO-based training of MoE models.

\paragraph{Benefit of GSPO} While Routing Replay makes it possible for GRPO training in MoE models to converge successfully, its mechanism of reapplying routing configurations imposes extra demands on memory and communication resources, and may restrict the model’s effective capacity. By contrast, GSPO, as depicted in Figure~\ref{fig:results}, avoids reliance on Routing Replay altogether, enabling straightforward calculation of the importance ratios $s_i(\theta)$, supporting stable optimization and standard convergence behavior. The crucial observation is that GSPO concentrates on the sequence likelihood (i.e., $\pi_{\theta} (y_i | x)$), rather than being influenced by the per-token likelihood (i.e., $\pi_{\theta} (y_{i,t} | x, y_{i,<t})$). Because the MoE model consistently preserves its ability for language modeling, the sequence likelihood exhibits limited variability. Ultimately, GSPO directly addresses the volatility in expert activation characteristic of MoE models, rendering methods such as Routing Replay unnecessary. As a result, the approach not only streamlines and reinforces the robustness of the training regime, but also enables the model to exploit its complete capacity absent artificial limitations.

\subsection{Benefit of GSPO for RL Infrastructure}

Due to the varying numerical precision between different engines used for training (such as Megatron) and those used during inference (such as SGLang and vLLM), the common approach is to recalculate the likelihoods of sampled outputs using the training engine, specifically for the old policy $\pi_{\theta_\text{old}}$. Nevertheless, since GSPO relies exclusively on sequence-level likelihoods as opposed to token-level ones, it exhibits a higher degree of robustness to these precision mismatches. Consequently, GSPO enables the use of likelihood estimates directly obtained from the inference engine for optimization, obviating the requirement for likelihood recalculation via the training engine. This characteristic proves particularly advantageous in contexts like partial rollout, multi-turn RL, and training-inference separated systems.

\section{Conclusion}

Group Sequence Policy Optimization (GSPO) is introduced as a novel algorithm in the field of reinforcement learning, designed for the training of large language models. GSPO builds upon the framework of importance sampling, formulating importance ratios grounded in sequence likelihood and conducting sequence-level processes such as clipping, rewarding, and optimization. When compared with GRPO, GSPO offers markedly better stability during training, greater efficiency, and enhanced overall performance. It proves especially effective in large-scale reinforcement learning for MoE models, underpinning the significant advancements found in recent Qwen3 models. As a scalable methodological foundation, GSPO enables continued expansion of reinforcement learning efforts, with anticipation of further transformative progress in artificial intelligence.

\end{document}